\documentclass[12pt]{article}
\usepackage[a4paper, margin=2.5cm]{geometry}
\usepackage{setspace}
\usepackage{titling}
\usepackage[T1]{fontenc}
\usepackage[utf8]{inputenc}
\usepackage{amsmath}
\usepackage{enumitem}
\usepackage{listings}
\usepackage{xcolor}
\usepackage{graphicx}
\usepackage{float}

\lstset{
    language=C++,
    basicstyle=\ttfamily\small,
    keywordstyle=\bfseries,
    commentstyle=\itshape,
    numbers=left,
    numberstyle=\tiny,
    frame=single,
    breaklines=true,
    tabsize=4
}

\begin{document}

\begin{titlepage}
    \centering

    {\large Ministerul Educaţiei și Cercetării al Republicii Moldova\par}
    \vspace{0.5cm}
    {\large Universitatea Tehnică a Moldovei\par}
    \vspace{0.5cm}
    {\large Facultatea Calculatoare, Informatică și Microelectronică\par}

    \vspace{3cm}

    {\Large \textbf{Laboratory Work 3:}\par}
    \vspace{0.5cm}
    {\Large \textbf{Empirical Analysis of Graph Traversal Algorithms}\par}
    \vspace{0.3cm}
    {\Large \textbf{Depth First Search and Breadth First Search}\par}

    \vspace{8cm}

    \begin{flushleft}
        Elaborated:\\
        st. gr. FAF-241 Caț Chiril

        \vspace{1cm}

        Verified:\\
        asist. univ. Reițman Patricia
    \end{flushleft}

    \vfill

    Chişinău -- 2026

\end{titlepage}

\tableofcontents
\newpage

\section{Algorithm Analysis}

\subsection{Objective}

The objective of this laboratory work is to implement and analyze two fundamental graph traversal algorithms: Breadth First Search (BFS) and Depth First Search (DFS), and to compare their practical performance through empirical measurements.

\subsection{Tasks}

According to the task description :contentReference[oaicite:0]{index=0}, the laboratory work includes:

\begin{enumerate}
    \item Implementation of BFS and DFS algorithms;
    \item Defining the properties of input data used for analysis;
    \item Choosing appropriate performance metrics;
    \item Performing empirical analysis;
    \item Graphical visualization of results;
    \item Drawing conclusions.
\end{enumerate}

\subsection{Theoretical Notes}

Graph traversal algorithms are fundamental tools used to explore the structure of graphs. They allow visiting all reachable nodes starting from a given source node.

Two main approaches exist:

\begin{itemize}
    \item Breadth First Search (BFS) explores nodes level by level;
    \item Depth First Search (DFS) explores nodes by going as deep as possible before backtracking.
\end{itemize}

Both algorithms have a theoretical time complexity of:

\[
O(V + E)
\]

where \(V\) is the number of vertices and \(E\) is the number of edges.

Although their asymptotic complexity is identical, their traversal patterns differ significantly, which can lead to differences in practical performance.

\subsection{Introduction}

Graphs are widely used to model relationships between entities in many domains such as networking, social systems, and routing problems. Efficient traversal of graphs is essential for solving many computational problems.

Breadth First Search is commonly used when the shortest path in an unweighted graph is required, while Depth First Search is often used for tasks such as connectivity checking, cycle detection, and topological ordering.

In this laboratory work, both algorithms are implemented and evaluated experimentally on randomly generated connected graphs.

\subsection{Comparison Metric}

The main metric used is execution time \(T(n)\), measured using CPU cycle counters and converted to seconds.

Additionally, throughput is measured as:

\[
\text{throughput} = \frac{\text{nodes visited}}{\text{time}}
\]

This allows evaluating not only total runtime but also efficiency relative to graph size.

\subsection{Input Format}

Graphs are generated programmatically as connected undirected graphs.

The generation process ensures:

\begin{itemize}
    \item Connectivity (tree backbone);
    \item Additional random edges;
    \item Increasing graph size from 100 to 10000 nodes.
\end{itemize}

This approach provides controlled yet realistic input data for performance evaluation.

\section{Implementation}

The algorithms were implemented in C++ and tested on randomly generated graphs.

\subsection{Graph Generation and Benchmarking}

\begin{lstlisting}
Graph generate_connected_graph(int n, int extra_edges) {
    std::mt19937 rng(std::random_device{}());
    std::uniform_int_distribution<int> dist(0, n - 1);

    Graph g(n);

    if (n <= 1)
        return g;

    for (int i = 1; i < n; ++i) {
        std::uniform_int_distribution<int> parent_dist(0, i - 1);
        int parent = parent_dist(rng);

        g[i].push_back(parent);
        g[parent].push_back(i);
    }

    int added = 0;
    while (added < extra_edges) {
        int u = dist(rng);
        int v = dist(rng);

        if (u == v)
            continue;

        if (std::find(g[u].begin(), g[u].end(), v) != g[u].end())
            continue;

        g[u].push_back(v);
        g[v].push_back(u);

        ++added;
    }

    return g;
}
\end{lstlisting}

\subsection{Breadth First Search (BFS)}

\subsubsection*{Algorithm Description}

BFS explores the graph layer by layer starting from the initial node. It visits all neighbors of a node before moving to the next level.

\subsubsection*{Implementation}

\begin{lstlisting}
inline std::vector<int> bfs(const Graph& graph, int start_node, auto&& fn) {
    std::vector<int> ans = { start_node };
    std::vector<int> current = { start_node };
    std::vector<bool> visited(graph.size(), false);

    visited[start_node] = true;

    fn(start_node);
    while (!current.empty()) {
        std::vector<int> next;

        for (const auto& v : current) {
            for (const auto& u : graph[v]) {
                if (!visited[u]) {

                    fn(u);
                    next.push_back(u);
                    visited[u] = true;
                    ans.push_back(u);
                }
            }
        }

        current = std::move(next);
    }

    return ans;
}
\end{lstlisting}

\subsubsection*{Results}

\begin{figure}[H]
    \centering
    \includegraphics[width=0.8\textwidth]{bfs.png}
    \caption{Quick Sort empirical results}
\end{figure}

The BFS traversal processes nodes level by level, which leads to predictable memory access patterns and stable performance.

\subsection{Depth First Search (DFS)}

\subsubsection*{Algorithm Description}

DFS explores as deep as possible along each branch before backtracking.

\subsubsection*{Implementation}

\begin{lstlisting}
inline std::vector<int> dfs(const Graph& graph, int start_node, auto&& fn) {
    std::vector<int> ans;
    std::vector<bool> visited(graph.size(), false);

    auto do_dfs = [&](auto&& self, int node) {
        if (visited[node]) return;

        fn(node);
        visited[node] = true;
        ans.push_back(node);

        for (const auto v : graph[node]) {
            self(self, v);
        }
    };

    do_dfs(do_dfs, start_node);
    return ans;
}
\end{lstlisting}

\subsubsection*{Results}

\begin{figure}[H]
    \centering
    \includegraphics[width=0.8\textwidth]{dfs.png}
    \caption{Quick Sort empirical results}
\end{figure}

DFS follows a depth-oriented traversal pattern, which may lead to different memory behavior compared to BFS.

\subsection{Benchmarking}

\begin{lstlisting}
template<typename TraversalFn>
void benchmark_traversal(const std::string& name,
                         TraversalFn&& traversal_fn) {

    const uint64_t cycles_per_sec = calibrate_tsc();
    const int repeats = 50;

    std::ofstream out("../results/lab2/" + name + ".csv");
    out << "n,throughput,time\n";

    for (int n = 100; n <= 10000; n += 100) {
        Graph g = generate_connected_graph(n, n);

        uint64_t total_cycles = 0;
        size_t total_nodes_visited = 0;

        for (int r = 0; r < repeats; ++r) {
            unsigned aux;

            uint64_t start = __rdtscp(&aux);
            auto visited = traversal_fn(g, 0, [](int){});
            uint64_t end = __rdtscp(&aux);

            total_cycles += (end - start);
            total_nodes_visited += visited.size();
        }

        double avg_cycles = (double)total_cycles / repeats;
        double seconds = avg_cycles / cycles_per_sec;
        double throughput = (double)(total_nodes_visited / repeats) / seconds;

        out << n << "," << throughput << "," << seconds << "\n";
    }
}
\end{lstlisting}

\section{Conclusion}

Both BFS and DFS have the same theoretical complexity \(O(V + E)\), but their practical behavior differs due to traversal patterns and memory access.

BFS provides a structured level-by-level traversal, which can lead to more predictable performance. DFS, on the other hand, follows a recursive path, which may result in deeper call stacks and different cache behavior.

The empirical analysis shows that despite identical asymptotic complexity, implementation details and traversal order influence real-world performance. Therefore, choosing between BFS and DFS depends not only on theoretical considerations but also on the specific problem and constraints.

\end{document}
